% А2.3 — Логарифм: определение, ограничения, знак
% Старая методичка: стр. 61 (старое задание 6)
% Источники: методичка «Логарифмы» (определение, lg и ln, «Откуда берутся ограничения»),
% шпаргалка «Логарифм» (знак логарифма). Стиль — стиль-изложения.md

\begin{tema}{А2.3}
\opr{Логарифмом} числа $b>0$ по основанию $a$ ($a>0$, $a\ne1$) называется показатель
степени, в которую нужно возвести $a$, чтобы получить $b$:
\[
  \log_a b=c \iff a^c=b.
\]
Например, $\log_3 9=2$, так как $3^2=9$; \ $\log_4 2=\frac12$, так как $4^{\frac12}=2$;
\ $\log_{\frac12}16=-4$, так как $\left(\frac12\right)^{-4}=16$.

\opr{Следствия из определения} (везде $a>0$, $a\ne1$):
\begin{pfTablica}{ll}
Формула & Почему\\
$\log_a 1=0$ & $a^0=1$\\
$\log_a a=1$ & $a^1=a$\\
$\log_a a^c=c$ & $a^c$ получается из $a$ возведением в степень $c$\\
$a^{\log_a b}=b$, \ $b>0$ & $\log_a b$ — степень, дающая $b$\\
\end{pfTablica}
Последнее равенство называют \emph{основным логарифмическим тождеством}.
Например, $\log_5 125=\log_5 5^3=3$, \ $7^{\log_7 3}=3$.

\opr{Десятичный логарифм} — логарифм по основанию 10: \ $\lg b=\log_{10}b$.
Например, $\lg1000=3$, \ $\lg0{,}01=-2$.

\opr{Натуральный логарифм} — логарифм по основанию $e\approx2{,}718$: \ $\ln b=\log_e b$.
Например, $\ln e=1$, \ $\ln1=0$.

\opr{Ограничения.} Выражение $\log_a b$ имеет смысл только при
$a>0$, $a\ne1$, $b>0$:
\begin{enumerate}[itemsep=0pt, topsep=2pt]
  \item $a>0$: в дробную степень можно возводить только положительные числа (\tref{А2.2});
  \item $a\ne1$: единица в любой степени равна 1, поэтому, например, $\log_1 2$
        не существует;
  \item $b>0$: $b=a^c$ — степень положительного числа, а она всегда положительна.
\end{enumerate}
В уравнениях и неравенствах ограничения становятся условиями на переменную (\tref{А5.3}).

\begin{primer}
При каких $x$ имеет смысл выражение $\log_{x-1}(5-x)$?
Решение: выпишем все три ограничения.
Основание: $x-1>0$ и $x-1\ne1$, то есть $x>1$ и $x\ne2$.
Число под логарифмом: $5-x>0$, то есть $x<5$.
Ответ: $1<x<5$, $x\ne2$, то есть $x\in(1;2)\cup(2;5)$.
\end{primer}

\opr{Знак логарифма:}
\begin{itemize}[itemsep=0pt, topsep=2pt]
  \item $\log_a b>0$, если $a$ и $b$ по одну сторону от 1 (оба больше 1 или оба
        между 0 и 1);
  \item $\log_a b<0$, если по разные стороны;
  \item $\log_a 1=0$.
\end{itemize}
Коротко: знак $\log_a b$ совпадает со знаком $(a-1)(b-1)$.
Например, $\log_2 5>0$, \ $\log_{0{,}5}0{,}3>0$, \ $\log_3 0{,}2<0$.

\opr{Оценка логарифма.} Чтобы понять, между какими целыми числами лежит $\log_a b$,
найдите соседние степени $a$, между которыми лежит $b$.
При $a>1$: $8<10<16$, то есть $2^3<10<2^4$, значит $3<\log_2 10<4$.
При $0<a<1$ знаки неравенств меняются: $\left(\frac12\right)^{-3}<10<\left(\frac12\right)^{-4}$,
значит $-4<\log_{\frac12}10<-3$.

\begin{primer}
Определите знак выражения $\log_3 0{,}5\cdot\log_{0{,}2}0{,}1$.
Решение: $3>1$, а $0{,}5<1$ — по разные стороны от 1, значит $\log_3 0{,}5<0$.
$0{,}2<1$ и $0{,}1<1$ — по одну сторону, значит $\log_{0{,}2}0{,}1>0$.
Произведение отрицательного и положительного чисел отрицательно.
Ответ: выражение отрицательно.
\end{primer}

\begin{oshibka}
Считают, что логарифм не бывает отрицательным. Ограничение $b>0$ относится к числу
под логарифмом, а сам логарифм может быть любым: $\log_2\frac18=-3$.
\end{oshibka}
\end{tema}
